% Compile with pdfLaTeX.
\documentclass[10pt,a4paper]{article}

\usepackage[T1]{fontenc}
\usepackage[utf8]{inputenc}
% Palatino and Helvetica are the pdfLaTeX counterparts closest to the
% TeX Gyre Pagella / Adventor pairing used in the original CV.
\usepackage{mathpazo}
\usepackage[scaled=0.92]{helvet}

\usepackage[a4paper,left=17mm,right=17mm,top=14mm,bottom=14mm]{geometry}
\usepackage{xcolor}
\usepackage{enumitem}
\usepackage{tabularx}
\usepackage{titlesec}
\usepackage{hyperref}
\usepackage{microtype}
\usepackage{needspace}

\definecolor{accent}{RGB}{35,132,158}
\definecolor{textgray}{RGB}{82,86,89}
\definecolor{lightgray}{RGB}{150,153,155}

\hypersetup{
  colorlinks=true,
  urlcolor=accent,
  linkcolor=accent,
  pdfauthor={Nikolay Tulupov},
  pdftitle={Nikolay Tulupov -- Reinforcement Learning and Control Systems Engineer}
}

\pagestyle{empty}
\setlength{\parindent}{0pt}
\setlength{\parskip}{0pt}
\setlength{\tabcolsep}{0pt}
\linespread{1.08}
\emergencystretch=2em

\titleformat{\section}
  {\needspace{4\baselineskip}\sffamily\bfseries\large\color{accent}}
  {}{0pt}{\MakeUppercase}
  [\vspace{0.10em}\color{accent}\titlerule]
\titlespacing*{\section}{0pt}{1.00em}{0.55em}

\setlist[itemize]{
  leftmargin=1.25em,
  itemsep=0.28em,
  topsep=0.32em,
  parsep=0pt,
  partopsep=0pt
}

\newcommand{\cvrole}[4]{%
  \begin{tabularx}{\textwidth}{@{}>{\raggedright\arraybackslash}X
    >{\raggedleft\arraybackslash}p{0.27\textwidth}@{}}
    {\sffamily\bfseries #1} & {\small\color{textgray}#4} \\
    {\itshape #2} & {\small\color{lightgray}#3}
  \end{tabularx}
  \vspace{0.12em}
}

\newcommand{\project}[2]{%
  \needspace{3\baselineskip}
  {\sffamily\bfseries #1}\par
  \vspace{0.12em}
  #2
}

\newcommand{\skillrow}[2]{%
  \textbf{#1}\hspace{0.45em}#2\par
  \vspace{0.13em}
}

\begin{document}

% Header
\begin{center}
  {\fontsize{24}{27}\selectfont\sffamily\bfseries\color{accent}
   NIKOLAY TULUPOV}\par
  \vspace{0.22em}
  {\large\sffamily\color{textgray}
   Reinforcement Learning \& Control Systems Engineer}\par
  \vspace{0.48em}
  {\small
    Moscow, Russia
    \enspace\textcolor{lightgray}{|}\enspace
    \href{tel:+79307297552}{+7 930 729-75-52}
    \enspace\textcolor{lightgray}{|}\enspace
    \href{mailto:tulupov.nd@phystech.edu}{tulupov.nd@phystech.edu}
    \enspace\textcolor{lightgray}{|}\enspace
    \href{https://t.me/oncoming_lane}{Telegram}
  }\par
  \vspace{0.16em}
  {\small
    \href{https://github.com/oncominglane}{GitHub}
    \enspace\textcolor{lightgray}{|}\enspace
    \href{https://www.researchgate.net/profile/Nikolay-Tulupov}{ResearchGate}
    \enspace\textcolor{lightgray}{|}\enspace
    \href{https://orcid.org/0009-0004-2017-2288}{ORCID: 0009-0004-2017-2288}
  }
\end{center}

\section{Profile}

Control systems and reinforcement learning engineer with over two years of
engineering and applied R\&D experience. Specializes in learning-based control
of PMSM electric drives, simulation, experimental validation, and engineering
software development using Python, PyTorch, and MATLAB/Simulink. Co-author of
an IEEE EDM 2026 publication and two registered software systems.

\section{Core Expertise}

\skillrow{Machine Learning \& RL:}{Python, PyTorch, reinforcement learning, TD3,
actor--critic methods, reward design, experiment design}
\skillrow{Control \& Modeling:}{control systems, PMSM/IPMSM drives, FOC,
mathematical modeling, MATLAB/Simulink, Motor-CAD}
\skillrow{Engineering \& Tools:}{C/C++, Git, Linux, CAN, STM32/AT32, power
inverters, test-bench experiments}

\section{Relevant Experience}

\cvrole{Moscow Institute of Physics and Technology}
       {Research Engineer, Robotics and Electric Drive Laboratory}
       {Moscow, Russia}{Mar 2026 -- Present}
\begin{itemize}
  \item Research and develop PMSM electric-drive control methods, with a focus
        on reinforcement learning for closed-loop control of physical systems.
  \item Developed and evaluated a physics-informed TD3 approach for end-to-end
        PMSM torque control and compared it with a classical FOC baseline.
  \item Build software pipelines for test-bench data processing, experiment
        automation, quantitative evaluation, and visualization.
  \item Automate Motor-CAD calculations and electric-motor geometry optimization.
  \item Conduct experimental validation and prepare scientific publications and
        technical presentations on electric-drive control.
\end{itemize}

\vspace{0.35em}
\cvrole{AmpereMagnete}
       {Engineering Intern, Artificial Intelligence Group}
       {Troitsk, Moscow, Russia}{Mar 2025 -- Mar 2026}
\begin{itemize}
  \item Developed and experimentally evaluated PMSM/IPMSM control and
        calibration methods, including reinforcement-learning-based approaches.
  \item Developed embedded firmware for an AT32-based power inverter and
        software for remote inverter control and test-bench operation.
  \item Implemented CAN-based telemetry acquisition, visualization, logging, and
        experimental data analysis.
  \item Developed OptiRunMotorCAD for automated calculations and motor-geometry
        optimization; contributed to the registered StandControl application.
\end{itemize}

\section{Education}

\cvrole{Moscow Institute of Physics and Technology}
       {M.Sc. in Applied Mathematics and Physics}
       {Moscow, Russia}{2026 -- 2028 (expected)}
{\small\textit{Specialization: E-Mobility -- Technologies, Materials and Control}}\par
\vspace{0.25em}
\cvrole{Moscow Institute of Physics and Technology}
       {B.Sc. in Applied Mathematics and Physics}
       {Moscow, Russia}{2022 -- 2026}
{\small\textit{Concentration: Radio Engineering and Computer Technology}}\par

\newpage

\section{Additional Engineering Experience}

\cvrole{TIENDA LLC}
       {Engineer, Technical Department}
       {Moscow, Russia}{Apr 2026 -- Present}
\begin{itemize}
  \item Continue development of a remote-control system for the Motorola XTL2500
        radio following the project's organizational transfer from PromSvyazRadio.
\end{itemize}

\vspace{0.3em}
\cvrole{PromSvyazRadio LLC}
       {Assistant to the Chief Engineer, Technical Department}
       {Moscow, Russia}{Feb 2024 -- Apr 2026}
\begin{itemize}
  \item Developed a distributed client--server system on Linux-based single-board
        computers for remote radio control, real-time bidirectional audio, and
        command processing.
  \item Implemented SB9600 protocol communication, hardware interfaces, and
        end-to-end hardware/software integration; contributed to the schematic,
        PCB, enclosure, testing, and technical documentation.
\end{itemize}

\section{Selected Research \& Engineering Projects}

\project{Physics-Informed TD3 for End-to-End Torque Control of a PMSM Drive}{
  \begin{itemize}
    \item Designed the state and action spaces, reward function, curriculum, and
          evaluation pipeline for a continuous-control TD3 agent.
    \item Introduced physics-informed critic objectives based on monotonicity
          constraints, auxiliary next-state prediction, and switch-aware TD weighting.
    \item Achieved a torque MAE of 0.079~N\,m and a 98.9\% success rate in the
          final evaluation, with performance close to classical FOC.
    \item Evaluated robustness to measurement noise and analyzed real-time MCU
          deployment constraints. Published at IEEE EDM 2026.
  \end{itemize}
}

\vspace{0.25em}
\project{OptiRunMotorCAD}{
  Developed a software application for automating Motor-CAD calculations and
  electric-motor geometry optimization. Registered as a computer program in
  2025
  (\href{https://new.fips.ru/registers-doc-view/fips_servlet?DB=EVM\&DocNumber=2025693472\&TypeFile=html}{No.~2025693472}).
}

\vspace{0.4em}
\project{StandControl}{
  Co-developed a client--server application for remote power-inverter control,
  CAN telemetry monitoring, visualization, logging, and test-bench operation.
  Registered as a computer program in 2026
  (\href{https://new.fips.ru/registers-doc-view/fips_servlet?DB=EVM\&DocNumber=2026614166\&TypeFile=html}{No.~2026614166}).
}

\section{Selected Coursework \& Training}

\href{https://mipt.ru/institute/events/alternativnyy-kurs-obuchenie-s-podkrepleniem-vesna-2026}{\textbf{Reinforcement Learning}},
MIPT (Spring 2026)\hspace{0.8em}\textcolor{lightgray}{|}\hspace{0.8em}
\href{https://huggingface.co/learn/deep-rl-course}{\textbf{Deep Reinforcement Learning Course}},
Hugging Face

\section{Publications \& Recognition}

\begin{itemize}
  \item N. Tulupov and D. Kanurin,
        \href{https://doi.org/10.1109/EDM69524.2026.11632234}{``Physics-Informed TD3 for End-to-End Torque Control of a PMSM Drive,''}
        \textit{2026 IEEE 27th International Conference of Young Professionals in
        Electron Devices and Materials (EDM)}, 2026.
        DOI: 10.1109/EDM69524.2026.11632234.
  \item N. Tulupov, S. Makov, D. Kanurin, and A. Tishin,
        \href{https://www.elibrary.ru/item.asp?id=90362673}{``PMSM Electric Drive Calibration Method Using a Reinforcement Learning Algorithm,''}
        \textit{XII International Conference Engineering \& Telecommunications
        (En\&T 2025)}, MIPT, Moscow, 2025; abstract published in the proceedings.
  \item N. Tulupov, S. Makov, D. Kanurin, and A. Tishin,
        \href{https://www.elibrary.ru/item.asp?id=89026196}{``Calibration Method for a Permanent-Magnet Synchronous Motor (PMSM) Using Deep Reinforcement Learning,''}
        \textit{III International Youth Conference ``AI Technologies in Science
        and Education,''} Moscow State University, Moscow, 2025; abstract
        published in the proceedings. Best Presentation Award, Exact and
        Engineering Sciences section.
  \item N. Tulupov and D. Kanurin,
        \href{https://mipt.ru/upload/\%D0\%9E\%D0\%9A\%D0\%9F\%D0\%98/01\%20\%D0\%A4\%D0\%A0\%D0\%9A\%D0\%A2-2026.pdf}{``Hybrid Control Architecture for a PMSM Electric Drive Using Reinforcement Learning,''}
        \textit{68th All-Russian Scientific Conference of MIPT}, Moscow, 2026;
        abstract published in the proceedings.
\end{itemize}

\section{Languages}

Russian -- Native \hspace{1.4em}\textcolor{lightgray}{|}\hspace{1.4em}
English -- B2 (Upper-Intermediate)

\end{document}
